\documentclass{article}
\usepackage[turkish]{babel}
\usepackage[T1]{fontenc}
\usepackage[utf8]{inputenc}
\usepackage{amsmath}

\title{Fonksiyonel analize giri\c{s}}
\author{A.~Y{\i}lmaz}
\date{2026}

\begin{document}
\maketitle

\section{Giri\c{s}}

Fonksiyonel analiz, sonsuz boyutlu uzaylar{\i}n yap{\i}s{\i}n{\i} inceleyen
matematik dal{\i}d{\i}r.  Bu \c{c}al{\i}\c{s}mada Banach uzaylar{\i} ve
lineer operat\"orler \"uzerinde duruyoruz.

% Turkish I/i distinction: dotted I (\.{I}) vs dotless i ({\i})
\.{I}stanbul \"Universitesi'nde yap{\i}lan ara\c{s}t{\i}rma, Hilbert
uzaylar{\i}ndaki izometrileri konu almaktad{\i}r.  \.{I}lk sonu\c{c}lar
umut vericidir.

\section{Temel kavramlar}

Bir Banach uzay{\i} $(X, \|\cdot\|)$ tam normlu bir uzayd{\i}r.
\"Onemli \"ornekler aras{\i}nda $L^p$ uzaylar{\i} ve $C[a,b]$ uzay{\i}
yer al{\i}r.

\begin{theorem}[A\c{c}{\i}k G\"onderme Teoremi]
$T \colon X \to Y$ s\"urjektif bir s{\i}n{\i}rl{\i} lineer operat\"or ise,
$T$ a\c{c}{\i}k bir g\"ondermedir.
\end{theorem}

Bu teorem, s\"urekli bijektif lineer operat\"orlerin tersinin de s\"urekli
oldu\u{g}unu g\"osterir.  G\"unl\"uk uygulamalarda \c{c}ok \"onemlidir.

\section{Sonu\c{c}}

Fonksiyonel analiz, diferansiyel denklemler ve kuantum mekani\u{g}i gibi
alanlarda vazge\c{c}ilmez bir ara\c{c}t{\i}r.

\end{document}
